%% Section body, shared verbatim by ../paper/main.tex (preprint) and
%% ../submission/body.tex (Information Sciences submission).  The section heading and
%% label live in the including file, because the two versions number sections the same
%% way but the submission adds the journal's required unnumbered sections after them.
\subsection{Assumptions}

Five assumptions are used, always together. They are stated in full in \Cref{app:assumptions} and
summarised here. \emph{Regularity} (\Cref{as:reg}): $F$ is twice continuously differentiable with
bounded Hessian, $H=\nabla^2F(\ts)\succ0$, and $\nabla^2F$ is Lipschitz near $\ts$.
\emph{Expected smoothness and moments} (\Cref{as:mom}): the second moment of the score is controlled by
the excess objective, with a $2+2\eta$ margin, and $\theta\mapsto\E[gg^{\!\top}]$ is continuous at
$\ts$. \emph{Curvature form} (\Cref{as:curv}): $\nabla^2F(\theta)=\E[\alpha\Psi\Psi^{\!\top}]$ with
$\alpha\ge0$, a uniform moment bound, and an $L^2$-Lipschitz condition in $\theta$ near $\ts$. \emph{Identifiability} (\Cref{as:ident}): $\ts$ is
the only stationary point of $F$ and elsewhere the gradient points away from it, which holds whenever
$F$ is convex with a unique minimiser. \emph{Geometric mixing} (\Cref{as:mix}): the stream is strictly
stationary with $\alpha_{\mathrm{mix}}(k)\le K\cmix^{\,k}$ for some $\cmix\in(0,1)$, the score has
$4+\nu'$ moments, and $S\succ0$.

\Cref{as:reg,as:mom,as:curv} are those of the independent-data analysis
\citep{godichonbaggioni2025adaptive}, with \Cref{as:curv} strengthened to a local Lipschitz condition
because the curvature terms are now dependent. \Cref{as:mix} is the only genuinely new assumption, and
\Cref{app:asdisc} lists the processes that satisfy it. Davydov's covariance inequality
\citep{davydov1968convergence} converts it into the geometric bound on score autocovariances that we use
repeatedly,
\begin{equation}
\opnorm{\Gamma_k}\;\le\;C_\Gamma\,\rho^{\,|k|},
\qquad \rho=\cmix^{\,\nu/(2+\nu)}\in(0,1),
\label{eq:geoGamma}
\end{equation}
so that $S$ in \eqref{eq:S} converges absolutely.

\subsection{The curvature recursion}

The first result is where dependence bites. Under independence the terms accumulated in
\eqref{eq:curv} form, conditionally on the iterates, an independent array and standard strong laws
apply; under \Cref{as:mix} they are dependent, and the gapped design is what restores control,
because the subsequence indexed by $\mathcal{C}$ is itself stationary with mixing coefficients
$\alpha_{\mathrm{mix}}(jm)\le K\cmix^{jm}$.

\begin{lemma}[The curvature estimate is never ill-conditioned]
\label{lem:ridge}
Under \Cref{as:curv}, with $\iota_n=c_\iota\varsigma_\star n^{-\qr}$, $\varsigma_\star>0$
fixed and $\qr\in(0,1)$, there is
a finite random $c_1>0$ with
\[
\lambda_{\min}(\Hb_t)\ \ge\ c_1\,n_t^{-\qr}
\qquad\text{and hence}\qquad
\opnorm{\Hb_t^{-1}}\ \le\ c_1^{-1}n_t^{\qr}=O(t^{\qr})
\]
for \emph{all} $t$, almost surely.
\end{lemma}

This is the whole purpose of the vanishing ridge, and it needs no assumption on the iterates: it
makes $\sum_t \gamma_t^2\opnorm{\Hb_{t-1}^{-1}}^2=\sum_tO(t^{2\qr-2})<\infty$, the square-summability
half of the Robbins--Monro condition, for \emph{every} $t$ rather than eventually, and without any
eigenvalue computation.

\begin{theorem}[Curvature rate]
\label{thm:curv}
Under \Cref{as:reg,as:curv,as:mix}, suppose $\theta_t\to\ts$ almost surely with
$\norm{\theta_t-\ts}=O(t^{-1/2}\log^{c}t)$ a.s.\ for some $c<\infty$. Then for every
$\nu_H<\qr$,
\[
\opnorm{\Hb_t-H}=O\!\left(t^{-\nu_H}\right)
\quad\text{and}\quad
\opnorm{\Hb_t^{-1}-H^{-1}}=O\!\left(t^{-\nu_H}\right)
\qquad\text{a.s.}
\]
\end{theorem}

\Cref{thm:curv} presupposes a rate for the iterates while the results below use
\Cref{thm:curv}, so the logical order matters. It is a four-rung ladder, each rung using only the
ones below it: \Cref{lem:ridge} from the ridge alone, with no assumption on the iterates; then
$\theta_t\to\ts$ almost surely (\Cref{thm:as}); then a preliminary rate (\Cref{lem:prelim}), which
is the one place the restriction $\qr<1/4$ is used; then \Cref{thm:curv}, and after it
\Cref{thm:rate,thm:clt,thm:l2,thm:lrv}. \Cref{app:ladder} states the ladder in full and
\Cref{app:proofs} is organised in that order.

\subsection{Consistency, rate, and the central limit theorem}

The independent-data proofs use that $\eps_t$ is a martingale difference with respect to
$\Fil_{t-1}=\sigma(\xi_1,\dots,\xi_{(t-1)b})$. Under \Cref{as:mix} it is not, and no block length repairs
this, because the conditional bias of a block average is $O(1/b)$ and does not vanish at fixed $b$. Our
substitute is Gordin's martingale--coboundary decomposition
$\eps_t=m_{t+1}+z_t-z_{t+1}$ with $(m_t)$ a stationary square-integrable martingale difference and
$\E[m_1m_1^{\!\top}]=S/b$ (\Cref{lem:gordin}); every weighted sum $\sum_t\gamma_t\Hb_{t-1}^{-1}\eps_t$
then splits into a martingale that converges by \Cref{lem:ridge} and a telescoping coboundary that
converges because one rank-one Riccati update moves the preconditioner by only $O(1/t)$
(\Cref{lem:incr}). This is what makes the argument work with weights that are \emph{random and adapted}
--- unavoidable here, since they contain $\Hb_{t-1}^{-1}$ --- and it is why the long-run covariance $S$,
rather than the instantaneous $\Gamma_0$, is what appears in the limit: it is the second moment of the
martingale part (\Cref{app:gordin}).

\begin{theorem}[Almost sure convergence]
\label{thm:as}
Under \Cref{as:reg,as:mom,as:curv,as:ident,as:mix}, the iterates \eqref{eq:step} satisfy
$\theta_t\to\ts$ almost surely, for every fixed block length $b\ge1$ and gap $m\ge1$.
\end{theorem}

\begin{theorem}[Rate]
\label{thm:rate}
Under the assumptions of \Cref{thm:as}, $\norm{\theta_t-\ts}^2=O(\log(N_t)/N_t)$ almost
surely, where $N_t=bt$.
\end{theorem}

\begin{theorem}[Central limit theorem with the long-run sandwich]
\label{thm:clt}
Under \Cref{as:reg,as:mom,as:curv,as:ident,as:mix}, for each \emph{fixed} block length $b\ge1$,
\[
\sqrt{N_T}\,(\theta_T-\ts)\;\rightsquigarrow\;\mathcal{N}\!\left(0,\;H^{-1}SH^{-1}\right),
\qquad N_T=bT,
\]
with $S$ as in \eqref{eq:S}. The same limit holds for the weighted averaged variant. If instead
$b=b_N\to\infty$ with $b_N=o(\sqrt N)$, the same limit holds provided the random constants
$\Lambda$ of \Cref{lem:series} are bounded in probability uniformly in $N$ along the
sequence $b_N$.
\end{theorem}

\noindent
We separate the two cases because only the first is proved unconditionally, and the difference
matters: \Cref{thm:lrv} takes $b\asymp\log N$, so the interval we actually recommend lives in the
second case. \Cref{app:growingb} identifies the step of the proof that needs the uniformity, shows
why the naive route around it loses a factor $T^{1/2}$, and gives what we can say in its place ---
including that coverage is flat in $b$ over $b\in[10,320]$ at $N=200{,}000$
(\Cref{tab:ablation}(b)), which is the composition the two theorems need. \Cref{sec:limits} lists
this as an open point rather than burying it.



\begin{proposition}[Blocking is free]
\label{prop:blockfree}
Under the assumptions of \Cref{thm:clt},
\[
\theta_T-\ts=-\frac{1}{N_T}H^{-1}\sum_{i\le N_T}g_i+R_T,
\qquad \norm{R_T}=o_p(N_T^{-1/2}),
\]
so the asymptotic variance in \Cref{thm:clt} is the same for every block length $b$.
\end{proposition}

\Cref{prop:blockfree} matters for the design argument: averaging the gradient over a block costs
nothing statistically, which is why we are free to choose $b$ for the covariance estimator. It also
answers the objection that single-sample streaming updates are more sample-efficient than batched ones
\citep{li2026stream} --- true for a step size $\gamma_t\asymp t^{-a}$ with $a<1$, and false here.

\subsection{The variance-inflation factor and the exact failure of dependence-blind intervals}

\begin{proposition}[Coverage of dependence-blind intervals]
\label{prop:coverage}
Let
\[
\kappa_j=\left[H^{-1}SH^{-1}\right]_{jj}\Big/\left[H^{-1}\Gamma_0H^{-1}\right]_{jj},
\]
and let $\Gz\to_p\Gamma_0$. Then the interval
$\theta_{T,j}\pm z_{\alpha/2}\{[\Hb^{-1}\Gz\Hb^{-1}]_{jj}/N\}^{1/2}$ has asymptotic coverage
\[
2\,\Phi\!\left(z_{\alpha/2}/\sqrt{\kappa_j}\right)-1 .
\]
\end{proposition}

The proof is one line and the content is entirely in $\kappa_j$: dependence-blind intervals do not
merely lose a little coverage, their coverage converges to a computable number below the nominal level.
For the homogeneous autoregressive design of \Cref{sec:exp} --- covariates $x_t=\phi x_{t-1}+\cdots$,
errors $u_t=\psi u_{t-1}+\cdots$, $x\perp u$ --- textbook geometric summation gives
\begin{equation}
\kappa_j=\frac{1+\phi\psi}{1-\phi\psi}\quad\text{for every }j,
\label{eq:kappa}
\end{equation}
so at $\phi=\psi=0.7$ a nominal $95\%$ interval covers $74.9\%$. We present \eqref{eq:kappa} only
because it makes \Cref{prop:coverage} checkable with no free parameters (\Cref{fig:coverage}); the
qualitative failure has been observed before \citep{liu2023mixing}, whose simulated coverage ``hovers
around $85\%$'' at $\kappa=2$, against our prediction of $83.4\%$.
